\chapter{能量相干如何在自由飞行中变成到达时间结构}
\label{chap:feqo-propagation}

相互作用刚结束时，能谱只告诉我们每条边带有多少人口。两条边带之间的相位也在电子态中，却不会出现在单次能量测量里。自由飞行不改变各边带人口，但不同能量有不同动量，它们到达同一平面时积累不同空间相位；这些相位重新叠加后，电子的到达时间可以出现比光学周期短得多的结构。

\section{固定距离传播的精确相位}
\label{chap:feqo-free-propagation}

一维向前运动的正能电子满足 \(E(p)=\sqrt{m^2c^4+c^2p^2}\)。反解正动量分支，\(p(E)=c^{-1}\sqrt{E^2-m^2c^4}\)。设相互作用后第 \(\ell\) 条边带能量为 \(E_\ell=E_0+\ell\hbar\omega>mc^2\)，振幅为 \(c_\ell\)，且只含向前传播支。下面的相位使用精确色散；将探测面上的各边带写成共同幅度归一化，则还要求 \(\max_\ell|v_\ell-v_0|/v_0\ll1\)，其中 \(v_\ell=\partial_pE[p(E_\ell)]\)。若速度差不可忽略，须保留能量归一化到通量归一化的 \(v_\ell^{-1/2}\) 因子，并以概率流而非单纯 \(|\Psi|^2\) 计算到达时间。在这个窄带条件下，位置 \(z=L\) 的波函数为
\begin{equation}
\Psi(L,t)=\sum_\ell c_\ell
\exp\!\left[\frac{\ii}{\hbar}p(E_\ell)L
-\frac{\ii}{\hbar}E_\ell t\right].
\label{eq:feqo-exact-drift-phase}
\end{equation}
每项模长保持 \(|c_\ell|\)，所以能谱不变；干涉项含 \(p(E_\ell)-p(E_{\ell'})\)，所以时间分布会变。

令 \(v_0=\partial_pE(p_0)\)、\(\tau=t-L/v_0\)，提取共同载波 \(e^{\ii p_0L/\hbar-\ii E_0t/\hbar}\)。留下的包络为
\begin{equation}
\psi(\tau;L)=\sum_\ell c_\ell
e^{-\ii\ell\omega\tau+\ii\Phi_\ell(L)},\qquad
\Phi_\ell(L)=\frac{L}{\hbar}[p(E_\ell)-p(E_0)]-\frac{\ell\omega L}{v_0}.
\label{eq:feqo-drift-envelope}
\end{equation}
这是按距离而非按共同传播时间比较相位；二者若混用，会把群延迟重复计算。

Taylor 定理给 \(p'(E_0)=1/v_0\) 和 \(p''(E_0)=-1/(\gamma_0^3mv_0^3)\)。故 \(\Phi_0=0\)、线性项严格抵消，
\begin{equation}
\Phi_\ell(L)=-\frac{\hbar\omega^2L}{2\gamma_0^3mv_0^3}\ell^2
+R_{3,\ell}^{(p)},\qquad
|R_{3,\ell}^{(p)}|\leq\frac{L|\ell\hbar\omega|^3}{6\hbar}
\max_{E\in\mathcal B}|p'''(E)|.
\label{eq:feqo-sideband-phase-expansion}
\end{equation}
\(\mathcal B\) 是实际占据的能量区间。整体相位不改概率，线性相位只平移时间原点，二次相位才能改变峰形。二阶近似的控制量是 \(\max|R_{3,\ell}^{(p)}|\ll1\)，传播距离越长越需检查它；窄能散本身不是充分理由。

若初始波包还有连续能量宽度 \(\sigma_E\)，每一条边带都带着一个邻近能量分布。上述离散式是这些峰足够窄且能量标签可分离时的近似；严格计算应把 \(\sum_\ell\) 外再加上能量积分，并先保持同一电子内部的相干叠加。

这一积分在高斯例子里可以做完。令入射能量相对 \(E_0\) 的偏移为 \(\epsilon\)，取已归一化的能量振幅
\(f(\epsilon)=(2\pi\sigma_E^2)^{-1/4}e^{-\epsilon^2/(4\sigma_E^2)}\)，从而 \(\int|f(\epsilon)|^2\dd\epsilon=1\)。假设 \(\sigma_E\ll\hbar\omega\)，不同 \(\ell\) 的能量峰可分开标号，则探测面振幅在共同的通量归一化因子之外为
\[
\Psi(L,t)=\sum_\ell c_\ell\int_{-\infty}^{\infty}\dd\epsilon\,
f(\epsilon)
e^{\ii p(E_\ell+\epsilon)L/\hbar-\ii(E_\ell+\epsilon)t/\hbar}.
\]
在前述各 \(v_\ell\) 相近、通量因子可共用的范围内，对每个 \(\ell\) 保留 \(p(E_\ell+\epsilon)=p(E_\ell)+\epsilon/v_\ell+O(\epsilon^2)\)，并要求 \(L\max_\ell|p''(E_\ell)|\sigma_E^2/\hbar\ll1\)。剩下的 Fourier 高斯积分给第 \(\ell\) 条边带的时间振幅正比于
\[
c_\ell e^{\ii p(E_\ell)L/\hbar-\ii E_\ell t/\hbar}
\exp\!\left[-\frac{\sigma_E^2}{\hbar^2}
\left(t-\frac{L}{v_\ell}\right)^2\right].
\]
单条边带的时间强度标准差是 \(\hbar/(2\sigma_E)\)；不同边带的中心又在 \(L/v_\ell\) 到达。若这些中心相差足以分开，即使能量振幅 \(c_\ell\) 原本相干，它们在同一时刻也不能充分重叠，理想的周期聚束图像就失效。若 \(L|p''|\sigma_E^2/\hbar\) 不小，上面的高斯还会因群速度色散展宽，应返回未展开的能量积分。

\begin{exercise}
从 \(p(E)=c^{-1}\sqrt{E^2-m^2c^4}\) 直接求 \(p'(E_0),p''(E_0)\)，核对\cref{eq:feqo-sideband-phase-expansion} 的量纲。
\end{exercise}


先让一组已知的边带振幅传播，再计算它们在时间轴上的干涉，便能把“聚束”变成一个可检验的概率密度，而非只画一条示意曲线。

\section{从边带和式算出阿秒时间密度}
\label{chap:feqo-attosecond}

光学周期 \(T_\omega=2\pi/\omega\) 内的归一化时间密度定义为 \(I(\tau;L)=|\psi(\tau;L)|^2/T_\omega\)，其中 \(\psi\) 取式~\eqref{eq:feqo-drift-envelope} 且 \(\sum|c_\ell|^2=1\)。这里把绝对到达时刻按模 \(T_\omega\) 折叠；它是一个周期内的相位或到达时间分布，不是宣称无限长的电子波列在整个实时间轴上的概率积分为一。若要预测有限脉冲的绝对到达分布，必须保留前节的连续能量包络。展开模平方，
\begin{equation}
I(\tau;L)=\frac1{T_\omega}\sum_{\ell,\ell'}c_\ell c_{\ell'}^*
e^{-\ii(\ell-\ell')\omega\tau}
e^{\ii(\Phi_\ell-\Phi_{\ell'})}.
\label{eq:feqo-temporal-density}
\end{equation}
周期积分中的 \(e^{-\ii(\ell-\ell')\omega\tau}\) 正交，故 \(\int_0^{T_\omega}I\dd\tau=\sum|c_\ell|^2=1\)。这项检查说明时间峰只是概率在一个周期内重新分配，没有创造电子。

要把能谱与时间图形的差别定量说清，按整数 \(m\) 收集式~\eqref{eq:feqo-temporal-density} 中频率为 \(m\omega\) 的项，定义
\begin{equation}
b_m(L):=\sum_{\ell\in\mathbb Z}c_\ell c_{\ell-m}^*
e^{\ii[\Phi_\ell(L)-\Phi_{\ell-m}(L)]},
\qquad
T_\omega I(\tau;L)=\sum_{m\in\mathbb Z}b_m(L)e^{-\ii m\omega\tau}.
\label{eq:feqo-bunching-harmonics}
\end{equation}
\(b_0=1\)、\(b_{-m}=b_m^*\)，故密度为实数且已归一。由柯西--施瓦茨不等式，\(|b_m|\leq(\sum|c_\ell|^2)^{1/2}(\sum|c_{\ell-m}|^2)^{1/2}=1\)。它是相邻 \(m\) 格的\emph{相干振幅}之和；只知道 \(P_\ell=|c_\ell|^2\) 不能求它。若任一相关边带相位被随机化，各项在重复实验中互相抵消，\(|b_m|\) 就变小。

对理想 PINEM 相位调制，交互区出口处的周期波函数是模长为一的 \(u(\theta)=e^{\beta e^{-\ii\theta}-\beta^*e^{\ii\theta}}\)。把 \(|u|^2=1\) 作 Fourier 展开，立即得到 \(b_m(0)=0\)（\(m\ne0\)）。因此边带在出口已经存在，时间密度却仍平坦；随后出现的时间尖峰不是光场直接把电子“切成脉冲”，而是自由色散使原先相消的跨边带项取得不同相位。

为看见峰怎样长出来，暂时只保留 \(\ell=-1,0,1\)，设 \(c_0=J_0(2b)\)、\(c_{+1}=J_1(2b)e^{\ii\phi}\)、\(c_{-1}=-J_1(2b)e^{-\ii\phi}\)，其中 \(b=|\beta|\)。在二次色散下令 \(\Phi_{\pm1}=\Phi_2/2\)、\(\Phi_0=0\)，则
\[
\psi_3=J_0(2b)-2\ii J_1(2b)e^{\ii\Phi_2/2}
\sin(\omega\tau-\phi),
\]
从而
\begin{align*}
T_\omega I_3={}&J_0^2(2b)+4J_1^2(2b)\sin^2(\omega\tau-\phi)\\
&+4J_0(2b)J_1(2b)\sin(\Phi_2/2)\sin(\omega\tau-\phi).
\end{align*}
若 \(\Phi_2=0\)，最后一项消失，弱耦合时密度变化从 \(b^2\) 才开始；传播到 \(|\Phi_2|\simeq\pi\) 时，一阶谐波的幅度达到本三边带模型的最大值。三边带截断的积分是 \(J_0^2+2J_1^2\)，与一的差来自舍弃的 \(|\ell|\geq2\) 人口，量级为 \(O(b^4)\)。强耦合时须恢复完整和式。

写 \(\Phi_\ell\simeq\tfrac12\ell^2\Phi_2\)，并由式~\eqref{eq:feqo-sideband-phase-expansion} 得 \(\Phi_2=-\hbar\omega^2L/(\gamma_0^3mv_0^3)\)。当 \(|\Phi_2|=2\pi\) 时，因 \(\ell^2\equiv\ell\pmod2\)，二次相位 \(e^{\ii\pi\ell^2}=e^{\ii\pi\ell}\) 只把时间图形平移半个周期。因此时间密度的 Talbot 复现尺度为
\begin{equation}
L_T=\frac{2\pi\gamma_0^3mv_0^3}{\hbar\omega^2}.
\label{eq:feqo-talbot-length}
\end{equation}
完整波函数不带时间平移地复现要到 \(2L_T\)。这项结论仍受三阶余项控制；若 \(\max|R_{3,\ell}^{(p)}|\gtrsim1\)，理想复现会被破坏。

初始能散、相位抖动与距离抖动要分别进入。若电子中心能量有宽度 \(\sigma_E\) 的窄高斯分布，不同电子的群延迟差一阶为 \(\delta\tau=-L\delta E/(\gamma_0^3mv_0^3)\)；能量较高的电子先到达。对到达时间作经典平均后，第 \(m\) 个时间谐波乘
\[
\exp\!\left[-\frac12\left(\frac{m\omega L\sigma_E}{\gamma_0^3mv_0^3}\right)^2\right].
\]
若每次制备的光学相位独立地以方差 \(\sigma_\phi^2\) 抖动，另一因子为 \(e^{-m^2\sigma_\phi^2/2}\)。距离分布则需对 \(\Phi_\ell(L)\) 再平均。这三种操作物理来源不同，不能用一个“模糊宽度”代替。

\begin{exercise}
对三边带公式积分一个周期，并计算 \(\Phi_2=0\) 与 \(\Phi_2=\pi\) 时的一阶时间谐波幅度。
\end{exercise}


真实传播中，电子还会经历未记录的扰动。它们不一定改变人口，却可能随机改变边带间的相位；于是上面的时间峰会变浅。

\section{环境带走相位信息时，哪一项先衰减}
\label{chap:feqo-decoherence}

电子的边带人口与边带相干是不同的矩阵元。设传播中电子边带 \(|\ell\rangle\) 分别把环境初态 \(|\varepsilon_0\rangle\) 变为 \(|\varepsilon_\ell(t)\rangle\)。对联合态求环境偏迹得到
\begin{equation}
\rho_{\ell\ell'}(t)=\rho_{\ell\ell'}(0)
e^{-\ii(E_\ell-E_{\ell'})t/\hbar}
\langle\varepsilon_{\ell'}(t)|\varepsilon_\ell(t)\rangle .
\label{eq:feqo-decoherence-factor}
\end{equation}
同一通道的环境态范数是一，所以人口可以不变；不同通道留下可区分的环境记录时，重叠模长小于一，时间密度\cref{eq:feqo-temporal-density} 中的相应拍频就变浅。这一计算已经解释了“退相干”，无需先背主方程名称。

可直接求解的电子模型是随机相位。令 \(L_e|\ell\rangle=\ell|\ell\rangle\)，噪声哈密顿量为 \(H_{\rm noise}(t)=\hbar\xi(t)L_e\)，其中实随机过程 \(\xi\) 单位为 \(\si{\per\second}\)，平均零，相关函数 \(C(t-s)=\mathbb E[\xi(t)\xi(s)]\)。在每次噪声实现中，\(\rho_{\ell\ell'}\) 多乘 \(e^{-\ii(\ell-\ell')X_t}\)，\(X_t=\int_0^t\xi(s)\dd s\)。若 \(\xi\) 为高斯过程，\(X_t\) 也是零均值高斯变量，特征函数 \(\mathbb E[e^{-\ii kX_t}]=e^{-k^2\mathbb E X_t^2/2}\) 给
\begin{equation}
D_{\ell\ell'}(t)=\exp\!\left[-\frac{(\ell-\ell')^2}{2}
\int_0^t\dd s\int_0^t\dd s'\,C(s-s')\right].
\label{eq:feqo-gaussian-dephasing}
\end{equation}
指数无量纲，\(D_{\ell\ell}=1\)，且高阶时间谐波因 \(|\ell-\ell'|\) 大而衰减更快。这一通道也可写成对随机酉算符 \(e^{-\ii X_tL_e}\) 的概率平均，因此保持正性和迹。

取 \(C(u)=\sigma^2e^{-|u|/\tau_c}\)。把正方形 \([0,t]^2\) 按 \(s>s'\) 对半，积分可手算为
\[
\int_0^t\dd s\int_0^t\dd s'\,C(s-s')
=2\sigma^2\bigl[t\tau_c-\tau_c^2(1-e^{-t/\tau_c})\bigr].
\]
故 \(t\ll\tau_c\) 时 \(D_{\ell\ell'}\simeq e^{-(\ell-\ell')^2\sigma^2t^2/2}\)，而 \(t\gg\tau_c\) 时衰减率趋向 \((\ell-\ell')^2\sigma^2\tau_c\)。只有相关时间 \(\tau_c\) 远短于观测时间及电子相干变化尺度，才可用常数速率的马尔可夫方程代替原来的双重积分。

在这个后一个时间窗内，令 \(D=\sigma^2\tau_c\)（单位 \(\si{\per\second}\)）。因为对任意矩阵元有 \([L_e,[L_e,\rho]]_{\ell\ell'}=(\ell-\ell')^2\rho_{\ell\ell'}\)，刚才的长时指数等价于
\begin{equation}
\frac{\dd\rho_e}{\dd t}
=-\frac{\ii}{\hbar}[H_0,\rho_e]
-D[L_e,[L_e,\rho_e]].
\label{eq:feqo-markov-dephasing}
\end{equation}
这里 \(H_0=\sum_\ell E_\ell|\ell\rangle\langle\ell|\)，并假设它与 \(L_e\) 对易。双重对易不改变对角元，却使第 \(m\) 阶拍频以速率 \(Dm^2\) 衰减。这个方程不是任意环境都满足的公理；它是相关时间短、噪声高斯且只耦合 \(L_e\) 时由上面的精确平均得到的极限。若环境能实际交换边带能量，还须增加改变人口的跃迁项。

若噪声只是每次电子入射前已固定、而实验未记录的初始能量，平均也会降低可见对比度，但它属于制备混合；\cref{eq:feqo-gaussian-dephasing} 属于传播中的随机耦合。若保存相关制备或环境记录，两者能够通过条件数据或回波操作区别。探测器分辨率则只在最终测量概率形成后作用，不能再算作这里的退相干率。

\begin{exercise}
把指数相关函数的双重积分化为 \(2\int_0^t(t-u)C(u)\dd u\)，完成积分并核对短时和长时两个极限。
\end{exercise}

